\songtitle{Lo que se ve}

\tag*{2026}\tag{folk-rock}\tag{spanish-lyrics}
\begin{notes}
Lyrics by Claude, written in an attempt to mimic Carlos Thompson's voice and songwriting style following specific prompts about the overall subject: a clear morning watched from a window over a Colombian street -- close to it all, without anyone asking for anything. Retrieved from Dvigitt's Suno page (V6-MINI, 18 October 2026).
\end{notes}
\begin{style}
Folk-rock, lightly produced and room-close with natural room tone; intimate close-mic baritone with Colombian pronunciation, stacked unprocessed harmonies on the hooks, spoken bridge and outro; warm, unhurried 84 BPM; fingerpicked acoustic guitar and soft brushed drums.
\end{style}
\begin{lyrics}
\annotation{Estrofa 1}\\
Entra el sol por la ventana antes de que me levante.\\
Desde aquí se ve la calle, un pedazo de cerro,\\
la señora que barre la acera por tercera vez,\\
el señor del tinto con su termo y su carrito.

\annotation{Estribillo}\\
Esto es lo que se ve\\
cuando el día está claro:\\
nada que no estuviera ayer,\\
y todo distinto bajo el sol.

\annotation{Estrofa 2}\\
Pasa un niño corriendo, con el uniforme mal puesto.\\
Pasa la buseta llena, con la puerta abierta y la música alta.\\
Un perro cruza sin mirar y todos frenan sin decir nada.\\
Nadie se queja. Todos saben cómo funciona esto.

\annotation{Estribillo}\\
Esto es lo que se ve\\
cuando el día está claro:\\
nada que no estuviera ayer,\\
y todo distinto con la luz.

\annotation{Estrofa 3}\\
Podría salir. El sol no cobra entrada.\\
Podría bajar, pedir un tinto, decir buenos días,\\
ser uno más en la acera, bajo la misma luz.\\
No es tristeza. Es una forma de estar cerca\\
sin que nadie me pida nada.

\annotation{Puente, hablado}\\
La luz llega a la mesa, al borde del vaso, a mi mano.\\
Está tibia. Eso es verdad, aunque no lo explique todo.\\
Hace años pensaba que mirar era perderse algo.\\
Hoy no sé. A veces mirar es lo que se puede.

\annotation{Estribillo final}\\
Esto es lo que se ve\\
cuando el día está claro:\\
nada que no estuviera ayer,\\
y yo, un poco más cerca de verlo.

\songpart{Outro, hablado, instrumentos fuera}
La señora termina de barrer.\\
El sol se mueve un palmo sobre la mesa.\\
Voy a salir.\\
Más tarde.
\end{lyrics}

